\section{Splitting mean recall on the released checkpoint, in full}\label{app:recall}

\paragraph{One pass, every model.} In PredCls the predictor's object probabilities are exact
one-hot vectors, so the codebase's other two counterfactual effects are not new models:
the Total Effect equals TDE, and the Natural Indirect Effect is identically zero. We therefore
ran the released causal MOTIFS checkpoint once, with \texttt{EFFECT\_TYPE TDE}, and patched
the predictor to write, for every ordered object pair of every test image, the visual
logits, the context logits, the context logits under the averaged context, the frequency
logits, and the logits it returned. The returned logits equal
$(\mathrm{vis}+\mathrm{ctx})+\mathrm{frq}-[(\mathrm{vis}+\mathrm{ctx}_{\mathrm{avg}})+\mathrm{frq}]$
with a largest absolute difference of exactly $0$ over all pairs, which is the cancellation
described in \cref{sec:sggaudit}. The baseline is $(\mathrm{vis}+\mathrm{ctx})+\mathrm{frq}$;
single branches and branch pairs are their own sums; the logit-adjusted baseline (LA$_\tau$)
subtracts $\tau\log\pi(y)$, with $\pi$ the predicate frequencies of the training split
($439{,}063$ relations).

\paragraph{Replaying the evaluator.} For model $m$, relation $i=(s,o,y_i)$ and any predicate
$y$, let $r^{(m)}_i(y)$ be the position at which the official evaluator would first match the
triplet $(s,o,y)$. Under the graph constraint each pair contributes its top non-background
predicate, pairs are ordered by that predicate's probability, and a predicted pair matches
$(s,o)$ when both objects carry the same classes and both boxes overlap with IoU at least
$0.5$; so $r^{(m)}_i(y)$ is the best rank among matching pairs whose top predicate is $y$.
Without the graph constraint every (pair, predicate) is ranked by its probability and
$r^{(m)}_i(y)$ is the best position of $(p,y)$ over matching pairs $p$. Relation $i$ is a
recall@$k$ hit under label $y$ iff $r^{(m)}_i(y)<k$, which makes recall an exact function of
the labelling. Matching through duplicate objects matters: a relation can be recalled by a
different pair of near-identical boxes, as the evaluator allows. A unit test checks the ranks
against the evaluator's own matching code on random images with duplicated objects,
repeated pairs and relabelled ground truth, for $k\in\{1,5,20,50,100\}$, and fails on each of
three deliberate faults. On the real run the re-scored metrics reproduce the official ones
(\cref{tab:recall-validation}). Three relations, all in one image with 35 objects, belong to
pairs the model never scored; they keep their exact ranks (they can only be matched through
duplicates) and are dropped from the probability-based analyses.

\begin{table}[h]
\centering\footnotesize
\caption{Official evaluator against the offline replay. The official TDE row is this run's
evaluation; the official baseline row is the archived evaluation of the same checkpoint.}
\label{tab:recall-validation}
% generated by experiments/make_cvpr_recall_tables.py -- do not edit
\begin{tabular}{@{}lrrrrr@{}}
\toprule
Model & R@50 & mR@50 & ng-R@50 & ng-mR@50 & mR@100 \\
\midrule
base, official & 0.6612 & 0.1459 & 0.8248 & 0.3260 & 0.1584 \\
base, re-scored & 0.6612 & 0.1459 & 0.8248 & 0.3260 & 0.1584 \\
TDE, official & 0.4588 & 0.2476 & 0.5687 & 0.2981 & 0.2870 \\
TDE, re-scored & 0.4588 & 0.2476 & 0.5687 & 0.2981 & 0.2870 \\
LA$_1$, re-scored & 0.6437 & 0.2343 & 0.8191 & 0.4173 & 0.2509 \\
LA$_{.5}$, re-scored & 0.6587 & 0.1845 & 0.8255 & 0.3712 & 0.1990 \\
\bottomrule
\end{tabular}

\end{table}

\paragraph{Micro-averaged mean recall.} The official metric averages each predicate's recall
over the images containing it, then over predicates, counting a predicate absent from the
test set as zero. Pooling relations instead gives
$mR=\sum_iw(y_i)\,h_i(y_i)$ with $w(y)=1/(K n_y)$, $h_i(y)=\ind\{r_i(y)<k\}$ and $n_y$ the
number of relations labelled $y$. With
$H_i(y)=N\,w(y)\,[h^{(1)}_i(y)-h^{(0)}_i(y)]$ the estimator of \cref{sec:method} applies
unchanged: \cref{thm:population,thm:sample} and the influence function use only the
antisymmetric pair structure, never properness, and because transport preserves each cell's
label counts it preserves every $n_y$. We restrict the pooling to identified relations, so
that $\hdT$ is exactly the change in micro-averaged mean recall over the relations the split
covers. It tracks the official change ($+0.0972$ against $+0.1017$ at $k=50$). Restricting
$w$ to one tier, or one predicate, splits that tier's or predicate's share; the parts add up
across tiers and predicates exactly.

\begin{table}[t]
\centering\footnotesize
\caption{Mean-recall split for every protocol and $k$. LA$_\tau$: logit-adjusted baseline;
ctx, vis+ctx: the baseline's context branch alone and without the frequency branch.}
\label{tab:recall-split}
% generated by experiments/make_cvpr_recall_tables.py -- do not edit
\begin{tabular}{@{}llrrrl@{}}
\toprule
Comparison & Protocol & $\Delta mR$ & group & case & case 95\% CI \\
\midrule
TDE vs base & mR@20 & $+.0598$ & $+.0462$ & $+.0135$ & {\scriptsize$[+.0101,+.0170]$} \\
 & mR@50 & $+.0972$ & $+.0841$ & $+.0130$ & {\scriptsize$[+.0093,+.0168]$} \\
 & mR@100 & $+.1280$ & $+.1099$ & $+.0182$ & {\scriptsize$[+.0144,+.0220]$} \\
 & ng-mR@20 & $-.0037$ & $-.0134$ & $+.0096$ & {\scriptsize$[+.0058,+.0135]$} \\
 & ng-mR@50 & $-.0337$ & $-.0464$ & $+.0127$ & {\scriptsize$[+.0080,+.0173]$} \\
 & ng-mR@100 & $-.0489$ & $-.0724$ & $+.0236$ & {\scriptsize$[+.0175,+.0296]$} \\
\midrule
LA$_{1}$ vs base & mR@20 & $+.0831$ & $+.0709$ & $+.0122$ & {\scriptsize$[+.0086,+.0158]$} \\
 & mR@50 & $+.0941$ & $+.0802$ & $+.0139$ & {\scriptsize$[+.0101,+.0176]$} \\
 & mR@100 & $+.0981$ & $+.0834$ & $+.0147$ & {\scriptsize$[+.0108,+.0187]$} \\
 & ng-mR@20 & $+.0754$ & $+.0671$ & $+.0084$ & {\scriptsize$[+.0041,+.0126]$} \\
 & ng-mR@50 & $+.0935$ & $+.0880$ & $+.0055$ & {\scriptsize$[-.0010,+.0120]$} \\
 & ng-mR@100 & $+.1061$ & $+.0973$ & $+.0088$ & {\scriptsize$[+.0006,+.0170]$} \\
\midrule
LA$_{.5}$ vs base & mR@20 & $+.0365$ & $+.0311$ & $+.0055$ & {\scriptsize$[+.0033,+.0077]$} \\
 & mR@50 & $+.0419$ & $+.0356$ & $+.0064$ & {\scriptsize$[+.0040,+.0088]$} \\
 & mR@100 & $+.0435$ & $+.0367$ & $+.0068$ & {\scriptsize$[+.0044,+.0092]$} \\
 & ng-mR@20 & $+.0376$ & $+.0319$ & $+.0057$ & {\scriptsize$[+.0022,+.0092]$} \\
 & ng-mR@50 & $+.0448$ & $+.0438$ & $+.0010$ & {\scriptsize$[-.0026,+.0047]$} \\
 & ng-mR@100 & $+.0514$ & $+.0485$ & $+.0029$ & {\scriptsize$[-.0033,+.0091]$} \\
\midrule
TDE vs LA$_{1}$ & mR@20 & $-.0234$ & $-.0247$ & $+.0013$ & {\scriptsize$[-.0024,+.0050]$} \\
 & mR@50 & $+.0031$ & $+.0039$ & $-.0008$ & {\scriptsize$[-.0050,+.0033]$} \\
 & mR@100 & $+.0300$ & $+.0265$ & $+.0034$ & {\scriptsize$[-.0010,+.0078]$} \\
 & ng-mR@20 & $-.0792$ & $-.0804$ & $+.0013$ & {\scriptsize$[-.0035,+.0060]$} \\
 & ng-mR@50 & $-.1272$ & $-.1344$ & $+.0072$ & {\scriptsize$[-.0001,+.0145]$} \\
 & ng-mR@100 & $-.1549$ & $-.1697$ & $+.0148$ & {\scriptsize$[+.0064,+.0232]$} \\
\midrule
ctx vs base & mR@20 & $+.0136$ & $+.0078$ & $+.0058$ & {\scriptsize$[+.0040,+.0076]$} \\
 & mR@50 & $+.0155$ & $+.0086$ & $+.0068$ & {\scriptsize$[+.0047,+.0089]$} \\
 & mR@100 & $+.0149$ & $+.0084$ & $+.0065$ & {\scriptsize$[+.0042,+.0087]$} \\
 & ng-mR@20 & $-.0014$ & $-.0072$ & $+.0058$ & {\scriptsize$[+.0030,+.0086]$} \\
 & ng-mR@50 & $-.0042$ & $-.0136$ & $+.0094$ & {\scriptsize$[+.0058,+.0129]$} \\
 & ng-mR@100 & $-.0040$ & $-.0170$ & $+.0130$ & {\scriptsize$[+.0089,+.0172]$} \\
\midrule
vis+ctx vs base & mR@20 & $+.0012$ & $-.0027$ & $+.0040$ & {\scriptsize$[+.0026,+.0053]$} \\
 & mR@50 & $-.0009$ & $-.0062$ & $+.0053$ & {\scriptsize$[+.0036,+.0069]$} \\
 & mR@100 & $-.0038$ & $-.0083$ & $+.0045$ & {\scriptsize$[+.0029,+.0060]$} \\
 & ng-mR@20 & $-.0060$ & $-.0124$ & $+.0065$ & {\scriptsize$[+.0040,+.0089]$} \\
 & ng-mR@50 & $-.0149$ & $-.0235$ & $+.0086$ & {\scriptsize$[+.0056,+.0117]$} \\
 & ng-mR@100 & $-.0197$ & $-.0327$ & $+.0129$ & {\scriptsize$[+.0091,+.0167]$} \\
\midrule
TDE vs ctx & mR@20 & $+.0462$ & $+.0384$ & $+.0077$ & {\scriptsize$[+.0044,+.0111]$} \\
 & mR@50 & $+.0817$ & $+.0755$ & $+.0062$ & {\scriptsize$[+.0025,+.0100]$} \\
 & mR@100 & $+.1132$ & $+.1015$ & $+.0117$ & {\scriptsize$[+.0080,+.0155]$} \\
 & ng-mR@20 & $-.0023$ & $-.0061$ & $+.0038$ & {\scriptsize$[+.0005,+.0072]$} \\
 & ng-mR@50 & $-.0295$ & $-.0328$ & $+.0033$ & {\scriptsize$[-.0008,+.0074]$} \\
 & ng-mR@100 & $-.0449$ & $-.0554$ & $+.0105$ & {\scriptsize$[+.0051,+.0160]$} \\
\bottomrule
\end{tabular}

\end{table}

\begin{table}[t]
\centering\footnotesize
\caption{The $mR@50$ split by tier (training instances: head $>10{,}000$, tail $<500$).}
\label{tab:recall-tiers}
% generated by experiments/make_cvpr_recall_tables.py -- do not edit
\begin{tabular}{@{}llrrrl@{}}
\toprule
Comparison & Tier & $\Delta mR@50$ & group & case & case 95\% CI \\
\midrule
TDE vs base & head & $-.0255$ & $-.0241$ & $-.0014$ & {\scriptsize$[-.0019,-.0009]$} \\
 & body & $+.1224$ & $+.1082$ & $+.0142$ & {\scriptsize$[+.0105,+.0180]$} \\
 & tail & $+.0003$ & $+.0001$ & $+.0002$ & {\scriptsize$[-.0001,+.0005]$} \\
\midrule
LA$_{1}$ vs base & head & $+.0009$ & $+.0006$ & $+.0003$ & {\scriptsize$[-.0000,+.0006]$} \\
 & body & $+.0932$ & $+.0796$ & $+.0136$ & {\scriptsize$[+.0099,+.0173]$} \\
 & tail & $+.0000$ & $+.0000$ & $+.0000$ & {\scriptsize$[+.0000,+.0000]$} \\
\midrule
LA$_{.5}$ vs base & head & $+.0013$ & $+.0011$ & $+.0002$ & {\scriptsize$[-.0001,+.0004]$} \\
 & body & $+.0407$ & $+.0345$ & $+.0062$ & {\scriptsize$[+.0039,+.0085]$} \\
 & tail & $+.0000$ & $+.0000$ & $+.0000$ & {\scriptsize$[+.0000,+.0000]$} \\
\midrule
TDE vs LA$_{1}$ & head & $-.0264$ & $-.0248$ & $-.0017$ & {\scriptsize$[-.0021,-.0012]$} \\
 & body & $+.0292$ & $+.0286$ & $+.0006$ & {\scriptsize$[-.0035,+.0048]$} \\
 & tail & $+.0003$ & $+.0001$ & $+.0002$ & {\scriptsize$[-.0001,+.0005]$} \\
\bottomrule
\end{tabular}

\end{table}

\begin{table}[t]
\centering\footnotesize
\caption{Within-cell AUC (comparison-weighted), image bootstrap with 200 replicates.}
\label{tab:recall-auc}
% generated by experiments/make_cvpr_recall_tables.py -- do not edit
\begin{tabular}{@{}lrrrl@{}}
\toprule
Comparison & AUC new & AUC old & $\Delta$AUC & 95\% CI \\
\midrule
TDE vs base & 0.5731 & 0.5705 & $+.0026$ & {\scriptsize$[-.0028,+.0086]$} \\
LA$_1$ vs base & 0.5705 & 0.5705 & $+.0001$ & {\scriptsize$[-.0000,+.0001]$} \\
LA$_{.5}$ vs base & 0.5705 & 0.5705 & $+.0000$ & {\scriptsize$[-.0000,+.0001]$} \\
TDE vs LA$_1$ & 0.5731 & 0.5705 & $+.0026$ & {\scriptsize$[-.0029,+.0086]$} \\
TDE vs ctx & 0.5731 & 0.5715 & $+.0016$ & {\scriptsize$[+.0008,+.0025]$} \\
vis+ctx vs base & 0.5705 & 0.5705 & $+.0001$ & {\scriptsize$[-.0000,+.0001]$} \\
ctx vs base & 0.5715 & 0.5705 & $+.0010$ & {\scriptsize$[-.0042,+.0068]$} \\
vis vs base & 0.5578 & 0.5705 & $-.0127$ & {\scriptsize$[-.0163,-.0090]$} \\
IETrans vs base & 0.5612 & 0.5705 & $-.0093$ & {\scriptsize$[-.0181,+.0004]$} \\
IETrans vs LA$_1$ & 0.5612 & 0.5705 & $-.0093$ & {\scriptsize$[-.0182,+.0002]$} \\
IETrans vs TDE & 0.5612 & 0.5731 & $-.0119$ & {\scriptsize$[-.0201,-.0022]$} \\
IETrans vs IETrans w/o prior & 0.5612 & 0.5612 & $+.0000$ & {\scriptsize$[+.0000,+.0000]$} \\
\bottomrule
\end{tabular}

\end{table}

\paragraph{Reading the case-level part of a thresholded metric.} Every change that adds a
constant to all log-odds within a cell --- logit adjustment at any $\tau$, or dropping the
frequency branch --- leaves the within-cell AUC where it was (\cref{tab:recall-auc}; the
residue of $10^{-4}$ is float32 rounding of the stored probabilities) and still shows a
significant case-level part of mean recall (\cref{tab:recall-split}): $+0.0064$ and $+0.0139$
for LA$_{.5}$ and LA$_1$, $+0.0053$ for removing the frequency branch. The reason is the
threshold. A relation is a hit when its own label wins a comparison, and moving the boundary
of that comparison flips the relations closest to it; those are the relations the model
already scored as most like the newly favoured label, and they carry it more often than the
others in their cell. The case-level part of a hit rate thus contains existing
discrimination spent at a new operating point. We read it against LA$_1$, whose group-level
change matches TDE's, and use the within-cell AUC as the direct test of new discrimination.
Without the graph constraint, where every (pair, predicate) competes for the top $k$
directly, TDE's case-level part against LA$_1$ is $+0.0072$, interval $[-0.0001,+0.0145]$.

\paragraph{The AUC.} For cell $c$ and labels $y\ne z$ present in it,
$\mathrm{AUC}_c(y,z)$ is the probability that a relation labelled $y$ has larger
$\log q(y)-\log q(z)$ than one labelled $z$, ties counting one half; the reported AUC
averages these with weight $n_c(y)n_c(z)$, the number of comparisons, which is the
$p(y)p(z)$ weighting of the pairwise form of $\dR_c$. A temperature divides every log-odds by
the same constant, and any change of the predicted distribution that is constant within a
cell adds the same constant to every log-odds in it; neither changes any within-cell order.
The statistic is computed exactly with ties in $O(M\log M)$ for $M$ compared relations;
bootstrap replicates re-weight images, so each costs one pass.

\paragraph{The proper-score split in every configuration.} The proper-score analyses use
the $183{,}639$ relations whose pair was scored. Top-1 predicate accuracy is $0.6981$ for
the baseline, $0.5577$ for TDE and $0.6820$ for LA$_1$; TDE also raises zero-shot recall to
$zR@50=0.1432$. \Cref{tab:proper-configs} lists the proper-score split of
\cref{sec:sggaudit} in every configuration we ran; the matched quadratic case-level part
of TDE against the baseline has a cyclic randomization $p=.555$. Raw outputs
give a significant case-level \emph{loss}; \cref{sec:simulation} showed that a calibration
difference alone manufactures one. Grouping by subject alone reverses the sign of the
group-level part (\cref{prop:coarsen}). Matched, the quadratic and log scores disagree on
TDE's case-level part against the baseline and against LA$_1$; LA$_1$, which moves no
within-cell ranking, has a positive case-level part under both, so the proper-score split
too responds to the operating point.

\begin{table}[t]
\centering\footnotesize
\caption{Proper-score split of TDE and of the logit-adjusted baseline (LA), every
configuration. Matched rows fit each model's temperature on a random half of the images and
audit the other half.}
\label{tab:proper-configs}
% generated by experiments/make_cvpr_recall_tables.py -- do not edit
\begin{tabular}{@{}lllrrrl@{}}
\toprule
Comparison & Outputs & Score & $\hdT$ & $\hdP$ & $\hdR$ & $\hdR$ 95\% CI \\
\midrule
TDE vs base & raw & quad. & $-0.11651$ & $-0.10296$ & $-0.01355$ & {\scriptsize$[-0.01504,-0.01206]$} \\
TDE vs base & raw & log & $+0.03277$ & $+0.15897$ & $-0.12619$ & {\scriptsize$[-0.13093,-0.12145]$} \\
TDE vs base & raw, subject only & quad. & $-0.11536$ & $+0.10175$ & $-0.21711$ & {\scriptsize$[-0.22124,-0.21298]$} \\
TDE vs base & raw, audit half & quad. & $-0.11548$ & $-0.10376$ & $-0.01172$ & {\scriptsize$[-0.01374,-0.00970]$} \\
TDE vs base & matched & quad. & $-0.18264$ & $-0.18258$ & $-0.00006$ & {\scriptsize$[-0.00120,+0.00109]$} \\
TDE vs base & matched & log & $-0.54760$ & $-0.59157$ & $+0.04396$ & {\scriptsize$[+0.04079,+0.04714]$} \\
LA$_1$ vs base & matched & quad. & $-0.02503$ & $-0.02895$ & $+0.00392$ & {\scriptsize$[+0.00334,+0.00450]$} \\
LA$_1$ vs base & matched & log & $-0.05095$ & $-0.07947$ & $+0.02852$ & {\scriptsize$[+0.02719,+0.02985]$} \\
TDE vs LA$_1$ & matched & quad. & $-0.15761$ & $-0.15363$ & $-0.00398$ & {\scriptsize$[-0.00510,-0.00285]$} \\
TDE vs LA$_1$ & matched & log & $-0.49665$ & $-0.51209$ & $+0.01544$ & {\scriptsize$[+0.01282,+0.01806]$} \\
LA$_{.5}$ vs base & matched & quad. & $-0.00358$ & $-0.00573$ & $+0.00215$ & {\scriptsize$[+0.00188,+0.00242]$} \\
LA$_{.5}$ vs base & matched & log & $-0.00209$ & $-0.01398$ & $+0.01189$ & {\scriptsize$[+0.01133,+0.01244]$} \\
IETrans vs base & matched & quad. & $-0.17695$ & $-0.16071$ & $-0.01623$ & {\scriptsize$[-0.01767,-0.01480]$} \\
IETrans vs base & matched & log & $-0.48127$ & $-0.42692$ & $-0.05435$ & {\scriptsize$[-0.05786,-0.05085]$} \\
IETrans vs LA$_1$ & matched & quad. & $-0.15191$ & $-0.13176$ & $-0.02015$ & {\scriptsize$[-0.02174,-0.01857]$} \\
IETrans vs LA$_1$ & matched & log & $-0.43032$ & $-0.34745$ & $-0.08287$ & {\scriptsize$[-0.08747,-0.07828]$} \\
IETrans vs TDE & matched & quad. & $+0.00569$ & $+0.02187$ & $-0.01618$ & {\scriptsize$[-0.01747,-0.01488]$} \\
IETrans vs TDE & matched & log & $+0.06633$ & $+0.16465$ & $-0.09832$ & {\scriptsize$[-0.10388,-0.09275]$} \\
\bottomrule
\end{tabular}

\end{table}

\subsection{TDE along its branches}\label{app:tdepath}

In PredCls the baseline's logits are $\mathrm{vis}+\mathrm{ctx}+\mathrm{frq}$ and TDE's are
$\mathrm{ctx}-\mathrm{ctx}_{\mathrm{avg}}$, so the change runs through two intermediate models
built from the same branches: the baseline without the frequency prior
($\mathrm{vis}+\mathrm{ctx}$), and the context term alone. The split is linear in the gain
vector, the AUC is a difference of per-model values, and in the matched proper-score split
each model keeps one held-out temperature throughout; so along this path every part adds up
to TDE's exactly, which the script asserts (\cref{tab:tde-path}). The order of the three
steps is a choice, and the steps are attributions along it, not order-free effects.

\begin{table}[h]
\centering\footnotesize
\caption{TDE's change from the baseline, step by step (graph-constrained $mR@50$ above;
no-graph-constraint $mR@50$ and the matched quadratic-score case-level part below). Each
column sums over the steps to TDE's row in \cref{tab:recall-split,tab:proper-configs}.}
\label{tab:tde-path}
% generated by experiments/make_cvpr_recall_tables.py -- do not edit
\begin{tabular}{@{}lrrrlrl@{}}
\toprule
Step & $\Delta mR@50$ & group & case & case 95\% CI & $\Delta$AUC & 95\% CI \\
\midrule
drop the frequency prior & $-.0009$ & $-.0062$ & $+.0053$ & {\scriptsize$[+.0036,+.0069]$} & $+.0001$ & {\scriptsize$[-.0000,+.0001]$} \\
drop the visual term & $+.0164$ & $+.0148$ & $+.0016$ & {\scriptsize$[-.0006,+.0037]$} & $+.0010$ & {\scriptsize$[-.0044,+.0068]$} \\
subtract the averaged context & $+.0817$ & $+.0755$ & $+.0062$ & {\scriptsize$[+.0025,+.0100]$} & $+.0016$ & {\scriptsize$[+.0008,+.0025]$} \\
\midrule
\multicolumn{7}{@{}l}{\emph{No graph constraint; last two columns: matched quadratic-score case-level part}} \\
drop the frequency prior & $-.0149$ & $-.0235$ & $+.0086$ & {\scriptsize$[+.0056,+.0117]$} & $+0.00915$ & {\scriptsize$[+0.00852,+0.00979]$} \\
drop the visual term & $+.0106$ & $+.0099$ & $+.0007$ & {\scriptsize$[-.0026,+.0040]$} & $-0.00178$ & {\scriptsize$[-0.00283,-0.00074]$} \\
subtract the averaged context & $-.0295$ & $-.0328$ & $+.0033$ & {\scriptsize$[-.0008,+.0074]$} & $-0.00743$ & {\scriptsize$[-0.00845,-0.00641]$} \\
\bottomrule
\end{tabular}

\end{table}

Two observations make the mechanism plain. Dropping the frequency prior, which is constant
within every cell, cannot move the within-cell AUC and does not; it moves mean recall by
$-0.0009$, but shifts the matched quadratic case-level part by $+0.00915$, another instance
of the operating-point effect. And the averaged-context term that TDE subtracts is almost a
constant: splitting each logit term's energy, after removing each row's mean, into the
vector shared by every relation, what the subject--object cell adds, and what varies within
cells (\cref{tab:logit-energy}), $99.96\%$ of $\mathrm{ctx}_{\mathrm{avg}}$ is the shared
vector, with a non-shared remainder of $0.021$ logits against a size of $1.10$. The shared
vector correlates $+0.75$ with the log training prior. Subtracting it is therefore, to first
order, a logit adjustment of the context branch by a learned prior; its small remainder is
what moves the within-cell AUC by $+0.0016$.

\begin{table}[h]
\centering\footnotesize
\caption{Energy of each logit term across the $183{,}639$ scored relations: shared by all
relations (global), set by the subject--object cell (between), varying within cells
(within); root-mean-square size of the term and of its non-global part; correlation of the
global vector with the log training prior $\log\pi$. The frequency term's within-cell
share is zero up to float32 rounding.}
\label{tab:logit-energy}
% generated by experiments/make_cvpr_recall_tables.py -- do not edit
\begin{tabular}{@{}lrrrrrr@{}}
\toprule
Logit term & rms & global & between & within & non-global rms & corr.\ with $\log\pi$ \\
\midrule
frq & 0.98 & 0.2976 & 0.7023 & 0.0001 & 0.820 & $+0.779$ \\
vis & 2.07 & 0.4875 & 0.4126 & 0.0999 & 1.485 & $+0.913$ \\
ctx(post) & 2.34 & 0.3891 & 0.5614 & 0.0495 & 1.830 & $+0.839$ \\
ctx(avg) & 1.10 & 0.9996 & 0.0001 & 0.0003 & 0.021 & $+0.754$ \\
TDE & 1.94 & 0.1045 & 0.8229 & 0.0726 & 1.831 & $+0.631$ \\
baseline & 5.07 & 0.4441 & 0.5178 & 0.0381 & 3.778 & $+0.878$ \\
\bottomrule
\end{tabular}

\end{table}

\subsection{The IETrans checkpoint}\label{app:ietrans}

\paragraph{What was run.} The authors release one Neural-Motifs PredCls checkpoint for
VG150 (their model zoo, Google Drive), saved at iteration $30{,}000$. It holds no
frequency-bias weights: the model was trained without one. Their evaluation command builds
the bias at test time from the training split's log predicate frequencies per object-class
pair, so the released model's logit is a learned context term plus a fixed per-cell
log-prior; we ran exactly that command, with the dump hook of \cref{app:recall} recording
both terms and their sum (equal to the returned logits with largest difference $0$). Their
annotation file matches the standard VG150 file array for array, except that $5{,}000$
training images were moved to a validation split; the test split is identical.

\paragraph{Test relations.} IETrans's data loader removes exact duplicate relations --- the
same subject object, object object and predicate annotated twice --- from the test set as
well as the training set, which leaves $152{,}226$ test relations where the standard protocol
keeps $183{,}642$. A duplicate shares its pair and its label with the relation that was kept,
so the evaluator matches it at the same ranks and the model scores it with the same
probabilities. We therefore mapped every standard relation to the IETrans row with the same
image, boxes, classes and predicate: the map is total, uses every IETrans row, and no key is
ambiguous. \Cref{tab:ietrans-protocols} reports IETrans on both test sets; the difference is
small, and every comparison in the paper uses the standard relations.

\begin{table}[h]
\centering\footnotesize
\caption{The released IETrans checkpoint under the official recall rules, on its own test
relations and on the standard ones.}
\label{tab:ietrans-protocols}
% generated by experiments/make_cvpr_recall_tables.py -- do not edit
\begin{tabular}{@{}lrrrrr@{}}
\toprule
Test relations & R@50 & mR@20 & mR@50 & mR@100 & ng-mR@50 \\
\midrule
IETrans's (152{,}226) & 0.4858 & 0.2884 & 0.3576 & 0.3905 & 0.5278 \\
standard (183{,}642) & 0.4891 & 0.2887 & 0.3587 & 0.3906 & 0.5291 \\
\bottomrule
\end{tabular}

\end{table}

\begin{table}[h]
\centering\footnotesize
\caption{Mean-recall split of the IETrans checkpoint against the MOTIFS baseline (with
tiers), the logit-adjusted baseline, TDE, and itself without the test-time prior. The
proper-score rows are in \cref{tab:proper-configs}, the AUC rows in \cref{tab:recall-auc}.}
\label{tab:ietrans-split}
% generated by experiments/make_cvpr_recall_tables.py -- do not edit
\begin{tabular}{@{}llrrrl@{}}
\toprule
Comparison & Protocol / tier & $\Delta mR$ & group & case & case 95\% CI \\
\midrule
IETrans vs base & mR@50 & $+.2204$ & $+.2143$ & $+.0060$ & {\scriptsize$[-.0031,+.0151]$} \\
 & ng-mR@50 & $+.2069$ & $+.2087$ & $-.0018$ & {\scriptsize$[-.0134,+.0099]$} \\
 & mR@50 head & $-.0169$ & $-.0132$ & $-.0038$ & {\scriptsize$[-.0042,-.0033]$} \\
 & mR@50 body & $+.1484$ & $+.1429$ & $+.0055$ & {\scriptsize$[+.0020,+.0090]$} \\
 & mR@50 tail & $+.0889$ & $+.0846$ & $+.0043$ & {\scriptsize$[-.0041,+.0127]$} \\
\midrule
IETrans vs LA$_{1}$ & mR@50 & $+.1263$ & $+.1341$ & $-.0079$ & {\scriptsize$[-.0173,+.0015]$} \\
 & ng-mR@50 & $+.1134$ & $+.1206$ & $-.0072$ & {\scriptsize$[-.0185,+.0040]$} \\
\midrule
IETrans vs TDE & mR@50 & $+.1232$ & $+.1302$ & $-.0070$ & {\scriptsize$[-.0164,+.0023]$} \\
 & ng-mR@50 & $+.2406$ & $+.2550$ & $-.0144$ & {\scriptsize$[-.0261,-.0028]$} \\
\midrule
IETrans vs IETrans w/o prior & mR@50 & $-.0133$ & $-.0046$ & $-.0088$ & {\scriptsize$[-.0171,-.0004]$} \\
 & ng-mR@50 & $-.0145$ & $-.0026$ & $-.0119$ & {\scriptsize$[-.0205,-.0032]$} \\
\bottomrule
\end{tabular}

\end{table}

\paragraph{Reading it.} Adding the test-time prior to IETrans's context term is a change
constant within every cell: it leaves the within-cell AUC exactly where it was, costs
$0.0133$ of mean recall, and its case-level part, $-0.0088$, is again an operating-point
effect. Against the baseline, the mean-recall gain is $97\%$ group-level and its case-level
part is not distinguishable from zero, while both matched proper scores give a case-level
loss (\cref{tab:proper-configs}). The comparison is between models of the same family trained
separately: IETrans released no plain baseline, and the MOTIFS-SUM baseline of
\cref{sec:sggaudit} differs from IETrans's Motifs predictor in its fusion of the visual
term.
